The crucial result concerning the solvability of equations is the inverse function theorem, which says that locally an equation is uniquely solvable if its linearization is.

\begin{Theorem}{Inverse function theorem}
\label{thm:banach-inverse-function}
Let $X$ and $Y$ be Banach spaces and let $U$ be an open neighborhood of the origin in $X$. Let $F:U\to Y$ be continuously differentiable and assume that $F^{\prime}(0)\in\mathcal{L}(X,Y)$ is one-to-one and onto. Then there exists an open neighborhood $V$ of $F(0)$ in $Y$ and a continuously differentiable mapping $G:V\to X$ with the property that $F(G(y))=y$. Moreover, $G(y)$ is the only sufficiently small solution $x$ of the equation $F(x)=y$.
\end{Theorem}

\begingroup
\renewcommand{\endProof}{\par\nobreak\hfill$\blacksquare$\par\vspace{0.1cm}}
\begin{Proof}
We consider the mapping $T_y$ defined by
\begin{equation}\label{eq:inverse-function-contraction-map-definition}
T_y(x)=F^{\prime}(0)^{-1}(y-F(0))-F^{\prime}(0)^{-1}(F(x)-F(0)-F^{\prime}(0)x).
\end{equation}
Let $M=\|F^{\prime}(0)^{-1}\|$. We claim: If $\epsilon$ is chosen small enough and $\|y-F(0)\|<\epsilon/(2M)$, then $T_y$ is a contraction mapping the closed $\epsilon$-ball in $X$ into itself. To see this, we note that
\begin{equation}\label{eq:inverse-function-contraction-ball-invariance-bound}
\|T_y(x)\|\leq M\|y-F(0)\|+M\|F(x)-F(0)-F^{\prime}(0)x\|.
\end{equation}
The first term on the right-hand side is less than $\epsilon/2$, and the second term is $o(\epsilon)$ by the definition of the derivative. Hence we find $\|T_y(x)\|<\epsilon$ if $\epsilon$ is sufficiently small. Moreover, note that
\begin{equation}\label{eq:inverse-function-contraction-map-derivative}
T_y^{\prime}(x)=-F^{\prime}(0)^{-1}(F^{\prime}(x)-F^{\prime}(0))
\end{equation}
has norm less than, say, $1/2$ if $\epsilon$ is small. The contraction property now follows from the equation
\begin{equation}\label{eq:inverse-function-contraction-integrated-derivative}
T_y(x)-T_y(z)=\int_0^1T_y^{\prime}(z+\theta(x-z))(x-z)\,\mathrm{d}\theta.
\end{equation}
By the Banach contraction principle (Theorem~\ref{thm:banach-contraction-fixed-point}), $T_y$ has a unique fixed point within the ball of radius $\epsilon$. We call this fixed point $G(y)$. It is easy to check that the equation $T_y(x)=x$ is equivalent to $F(x)=y$.

It remains to be shown that $G$ is in fact continuously differentiable. We first show continuity. We have $G(y)=T_y(G(y))$ and hence
\begin{equation}\label{eq:inverse-function-lipschitz-continuity-estimate}
\begin{aligned}
\|G(y+h)-G(y)\|&\leq\|T_{y+h}(G(y+h))-T_y(G(y+h))\|\\
&\quad+\|T_y(G(y+h))-T_y(G(y))\|\\
&\leq M\|h\|+\frac{1}{2}\|G(y+h)-G(y)\|,
\end{aligned}
\end{equation}
where the latter estimate follows from the contraction property of $T_y$. From \eqref{eq:inverse-function-lipschitz-continuity-estimate}, it is immediate that $G$ is Lipschitz continuous.

Let now $k=G(y+h)-G(y)$, then it is immediate from \eqref{eq:inverse-function-lipschitz-continuity-estimate} that $\|k\|\leq2M\|h\|$. We now have
\begin{equation}\label{eq:inverse-function-differentiability-remainder-expansion}
0=F(G(y+h))-F(G(y))-h=F^{\prime}(G(y))k+R(k)-h,
\end{equation}
where $\|R(k)\|/\|k\|\to0$ as $\|k\|\to0$. It follows that
\begin{equation}\label{eq:inverse-function-increment-linearization}
k=F^{\prime}(G(y))^{-1}h-F^{\prime}(G(y))^{-1}R(k),
\end{equation}
and since $\|k\|\leq2M\|h\|$, we have that $\|R(k)\|/\|h\|\to0$ as $\|h\|\to0$. This proves that $G$ is differentiable and the derivative is $F^{\prime}(G(y))^{-1}$.
\end{Proof}
\endgroup

\begin{Remark}
\label{rem:inverse-function-higher-differentiability}
If $F$ is $k$ times continuously differentiable, then so is $G$ (Problem~\ref{ex:inverse-function-higher-differentiability}).
\end{Remark}
